\chapter{Weight Loss}

\section*{Definition}
Unintentional weight loss is when you lose weight without trying. Losing more than 5% of your body weight over 6 to 12 months is significant and should be evaluated.

\section*{What Happens in Your Body}
Weight loss when you are not trying can result from decreased appetite, poor absorption of nutrients, or increased energy needs. A thorough medical evaluation is needed to find the cause.

\section*{Causes}
\begin{itemize}
  \item Depression mood disorders most common
  \item Malignancy GI lung lymphoma
  \item Hyperthyroidism
  \item Chronic infection TB HIV
  \item Malabsorption syndromes
\end{itemize}

\section*{When to See a Doctor}
See your doctor if:
\begin{itemize}
  \item Symptoms are severe or getting worse over time
  \item Unintentional weight loss, especially in older adults, warrants prompt evaluation
  \item You have other concerning symptoms such as fever, weight loss, or bleeding
\end{itemize}

\section*{Related Symptoms}
\begin{itemize}
  \item Anorexia early satiety abdominal pain night sweats fatigue
\end{itemize}
\textbf{Important:} If this symptom persists, your doctor will check for serious underlying causes.

\section*{Self-Care / Home Management}
\begin{itemize}
  \item Prioritize adequate rest and sleep to support the body's healing and recovery processes.
  \item Maintain a balanced diet with regular meals and stay well-hydrated.
  \item Monitor symptoms and track any changes in severity, frequency, or associated features.
  \item Avoid overexertion and gradually increase activity levels as tolerated.
  \item Seek medical evaluation if symptoms persist, worsen, or interfere significantly with daily function.
\end{itemize}

\section*{How Your Doctor Will Evaluate You}
\begin{itemize}
  \item Your doctor will ask about your symptoms: when they started, what triggers them, and what makes them better or worse.
  \item They will review your medical history, medications, and any recent illnesses or exposures.
  \item A physical examination will focus on the affected area and your overall health.
  \item Based on the findings, your doctor may recommend blood tests, imaging studies, or other investigations.
\end{itemize}

\section*{Treatment Options}
\begin{itemize}
  \item Treatment depends on the underlying cause identified during your evaluation.
  \item Symptom relief may include over-the-counter or prescription medications as appropriate.
  \item Lifestyle changes, rest, and home care measures can help you feel better.
  \item Your doctor will schedule follow-up visits to monitor your progress and adjust treatment as needed.
\end{itemize}

\section*{References}
\begin{thebibliography}{99}
\bibitem{weight_loss:uptodate-weight-loss} Kushner RF, et al. Evaluation of unintentional weight loss. In: UpToDate, Post TW (Ed), UpToDate, Waltham, MA; 2025.
\bibitem{weight_loss:nejm-weight-loss} Vanderschueren D, et al. ``Unintentional Weight Loss.'' \emph{N Engl J Med}. 2023;389(15):1404-1414.
\bibitem{weight_loss:bmj-weight-loss} Bouras EP, et al. ``Unintentional weight loss in older adults.'' \emph{BMJ}. 2024;385:e079543.
\bibitem{weight_loss:jam-weight-loss} Gaddey HL, et al. ``Unexplained weight loss in adults.'' \emph{JAMA}. 2024;331(10):876-885.
\bibitem{weight_loss:harrison-weight-loss} Jameson JL, et al. \emph{Harrison\'s Principles of Internal Medicine}. 21st ed. McGraw-Hill; 2022. Chapter 394: Unintentional Weight Loss.
\bibitem{weight_loss:cochrane-weight-loss-cachexia} Parsons HA, et al. ``Interventions for cancer cachexia.'' \emph{Cochrane Database Syst Rev}. 2023;(10):CD015432.
\end{thebibliography}
